 \documentclass{article}
\usepackage{amsmath}
\usepackage{graphicx}
\graphicspath{ {images/} }
\usepackage[spanish]{babel}
\usepackage{bbm}
\usepackage{amsthm}
\usepackage{authblk}
\usepackage{hyperref}
\usepackage[utf8]{inputenc}
\usepackage{mathrsfs}
\usepackage{amsfonts}
\usepackage{amssymb}
\usepackage{enumerate}
\usepackage[left=3cm,right=2cm,top=2cm,bottom=2cm]{geometry}
\usepackage{epsfig}
 \renewcommand{\baselinestretch}{0.93}
 \thispagestyle{empty}
 \pagestyle{empty}
\setlength{\textheight}{53pc} \setlength{\textwidth} {36pc}
\setlength{\topmargin}{-4mm}
\usepackage{multicol}
\newcommand*{\email}[1]{%
    \normalsize\href{mailto:#1}{#1}\par
    }
\setlength{\columnsep}{1cm}
\newcommand{\N}{\mathbb{N}}
\newcommand{\calM}{\cal{M}}
\newcommand{\calP}{\cal{P}}
\newcommand{\F}{\mathbb{F}}
\newcommand{\R}{\mathbb{R}}
\newcommand{\Z}{\mathbb{Z}}
\newcommand{\Q}{\mathbb{Q}}
\newcommand{\C}{\mathbb{C}}
\newcommand{\bbH}{\mathbb{H}}


\theoremstyle{definition}
\newtheorem{ejer}{Ejercicio}
\author{Marcos Morales}
\affil{Universidad Finis Terrae}
\title{Semana 4\\}



	


\begin{document}


\begin{picture}(400, 75)
   \put(-40,15){\includegraphics[width=5cm]{UFT.png}}


\end{picture}


\begin{center}
\huge{Laplace y Aplicaciones}
\end{center}

\begin{center}
\Large{ Marcos Morales, Camila Muñoz}
\end{center}

\begin{center}
	\large{Ecuaciones Diferenciales Ordinarias (EDO)}
\end{center}
\begin{center}
\setlength{\parindent}{0cm}\large{Clases centradas para capacitar a los estudiantes de ingeniería en Ecuaciones diferenciales ordinarias. \\}
\end{center}


\vspace{1cm}

\begin{center}
	\large{Contenidos:}
\end{center}

\setlength{\parindent}{2.95cm}- Segundo Teorema de Traslación\\
\phantom{abefwfewefwfefffwi} - Derivada de una Transformada de Laplace\\
\phantom{abefwfewefwfefffwi} - Integral de una transformada de Laplace\\
\phantom{abefwfewefwfefffwi} - Resolución de sistemas de ecuaciones diferenciales\\





\vspace{6cm}


\begin{center}
\today
\end{center}
\begin{center}
Santiago, Chile
\end{center}


\pagestyle{empty}


\setcounter{page}{1}
\pagestyle{plain}
\setlength{\parindent}{0cm}





\newpage
\section{Segundo teorema de traslación}



\textbf{Teorema (Segundo teorema de traslación): }Sea $f: [0,+\infty) \to \mathbb{R}$ continua por trozos y de orden exponencial. Si $a \in \mathbb{R}^+$, entonces

\begin{center}
$\mathcal{L}[u(t-a)f(t-a)](s) =e^{-as}\mathcal{L}[f(t)](s)$
\end{center}

O bien equivalentemente


\begin{center}
$\mathcal{L}[u(t-a)f(t)](s) =e^{-as}\mathcal{L}[f(t+a)](s)$
\end{center}

Por ejemplo $\displaystyle \mathcal{L}[u(t-2)\cos(t-2)](s) =e^{-2s}\mathcal{L}[\cos(t)](s) = \frac{e^{-2s} s}{s^2+1}$. \\ \\ \\



\textbf{Teorema (Segundo teorema de traslación(Versión transformada inversa)): }Sea $f: [0,+\infty) \to \mathbb{R}$ continua por trozos y de orden exponencial. Si $a \in \mathbb{R}^+$, entonces

\begin{center}
$\mathcal{L}^{-1}[e^{-as}f(s)](t) =u(t-a)\mathcal{L}^{-1}[f(s)](t-a)$
\end{center}


Por ejemplo $\displaystyle \mathcal{L}^{-1}\left[ \frac{e^{-2s}}{s^2} \right] (t)  = u(t-2) \mathcal{L}^{-1} \left[ \frac{1}{s^2} \right] (t-2) = (t-2)u(t-2)$ \\



\newpage


\begin{center}
\textbf{\large{Ejemplo con ecuaciones diferenciales}} \\
\end{center}

Vamos a resolver el siguiente problema de valores iniciales usando la transformada de Laplace


\begin{align*}
y''+4y &=f(t) \\
y(0)&=0 \\
y'(0)&=0
\end{align*}

Donde 

\begin{center}
$f(t) = \left\{ \begin{array}{lcc} 1 & si & 0 \leq t < 3 \\ \\ 2 & si & t \geq 3  \end{array} \right.$
\end{center}

\textit{Respuesta: } Aplicamos transformada de Laplace, tenemos que

\begin{align*}
\mathcal{L}[y''+4y] &=\mathcal{L}[f(t)] \\
\mathcal{L}[y'']+\mathcal{L}[4y] &=\mathcal{L}[f(t)]  \\
s^2\mathcal{L}[y]-sy(0)-y'(0)+4\mathcal{L}[y] &=\mathcal{L}[f(t)]\\
\mathcal{L}[y](s^2+4) &=\mathcal{L}[f(t)]
\end{align*}

Por otra parte, notemos que 

\begin{center}
$f(t ) = 1+(2-1)u(t-3)=1+u(t-3)$
\end{center}

Luego


\begin{center}
$\displaystyle \mathcal{L} [f(t)] = \mathcal{L}[1-u(t-3)] = \frac{1}{s}+\frac{e^{-3s}}{s} $
\end{center}

Luego



\begin{align*}
\mathcal{L}[y](s^2+4) &=\frac{1}{s}+\frac{e^{-3s}}{s}
\end{align*}


De donde se deduce que

\begin{align*}
\mathcal{L}[y] &=\frac{1}{s(s^2+4)}+\frac{e^{-3s}}{s(s^2+4)}
\end{align*}

Ahora usamos fraciones parciales

\begin{align*}
\frac{1}{s(s^2+4)} &= \frac{A}{s}+\frac{Bs+C}{s^2+4} \\
&= \frac{A(s^2+4)+s(Bs+C)}{s(s^2+4)} \\
&= \frac{s^2(A+B)+Cs+4A}{s(s^2+4)} \\
\end{align*}

De donde se deduce que $A=1/4$, $B=-1/4$ y $C=0$, luego


\begin{align*}
\frac{1}{s(s^2+4)} &= \frac{1}{4s}-\frac{s}{4(s^2+4)} \\
\end{align*}

Aplicando transformada inversa tenemos	


\begin{align*}
\mathcal{L}^{-1}\left[  \frac{1}{s(s^2+4)}\right]  &= \mathcal{L}^{-1}\left[\frac{1}{4s}-\frac{s}{4(s^2+4)} \right] \\
&= \frac{1}{4}-\frac{\cos(4t)}{4}
\end{align*}

Y aplicando el segundo teorema de traslación obtenemos que


\begin{align*}
\mathcal{L}^{-1}\left[  e^{-3s}\frac{1}{s(s^2+4)}\right]  &= u(t-3)\mathcal{L}^{-1}\left[\frac{1}{4s}-\frac{s}{4(s^2+4)} \right](t-3) \\
&= u(t-3)\left( \frac{1}{4}- \frac{\cos(4(t-3))}{4} \right) 
\end{align*}


Por lo tanto, la ecuación que teníamos

\begin{align*}
\mathcal{L}[y] &=\frac{1}{s(s^2+4)}+\frac{e^{-3s}}{s(s^2+4)}
\end{align*}

Se resuelve como

\begin{align*}
y &= \mathcal{L}^{-1} \left[ \frac{1}{s(s^2+4)}+\frac{e^{-3s}}{s(s^2+4)}\right] 
\end{align*}



\begin{align*}
y &= \frac{1}{4}-\frac{\cos(2t)}{4} +u(t-3)\left( \frac{1}{4}- \frac{\cos(2(t-3))}{4} \right) 
\end{align*}

O bien



\begin{center}
$y = \left\{ \begin{array}{lcc} \frac{1}{4}-\frac{\cos(2t)}{4} & si &  0\leq t < 3 \\ \\ \frac{1}{2}-\frac{\cos(2t)+\cos(2(t-3))}{4} & si & t \geq 3  \end{array} \right.$
\end{center}













\newpage




\section{Derivada de una transformada de Laplace}

El siguiente resultado puede ser demostrado directamente usando la definición de la transformada de Laplace, en este curso no vamos a hacer la demostración. \\

\textbf{Proposición: } Sea $f$ una función de orden exponencial y $n \in \mathbb{N}$, entonces

\begin{align*}
    (-1)^n \frac{d^n}{ds^n} \mathcal{L}[f](s) = \mathcal{L}[t^nf(t)](s)
\end{align*}

Este resultado nos sirve principalmente para calcular transformadas de Laplace de la forma $\mathcal{L}[t^nf(t)](s)$, para $n=1$ tenemos


\begin{align*}
     \mathcal{L}[tf(t)](s) = -\frac{d}{ds} \mathcal{L}[f](s)
\end{align*}

\textbf{Ejemplo: } Calculemos la transformada de Laplace de $f(t) =t\sin(at) $, sabemos que 

\begin{align*}
    \mathcal{L}[\sin(at)](s) = \frac{a}{s^2+a^2}
\end{align*}

Derivemos esto respecto a $s$, obtenemos

\begin{align*}
    \frac{d}{ds} \mathcal{L}[\sin(at)](s) &=  \frac{d}{ds} \left( \frac{a}{s^2+a^2} \right) \\
    &=  \frac{(s^2+a^2)a'-(s^2+a^2)'a}{(s^2+a^2)^2} \\
    &=  \frac{-2as}{(s^2+a^2)^2} \\
\end{align*}

Por lo tanto


\begin{align*}
    \mathcal{L}[t\sin(at)](s) = \frac{2as}{(s^2+a^2)^2} 
\end{align*}

Notemos que esto implica que

\begin{align*}
    \mathcal{L}^{-1} \left[\frac{s}{(s^2+a^2)^2} \right] = \frac{t\sin(at)}{2a} 
\end{align*}


\textbf{Ejercicio 1: } Demuestre que

\begin{align*}
    \mathcal{L}[t\cos(at)](s) = \frac{s^2-a^2}{(s^2+a^2)^2} 
\end{align*}



\textbf{Ejercicio 2: } Demuestre que

\begin{align*}
    \mathcal{L}[\sin(at)-at\cos(t)](s) = \frac{2a^3}{(s^2+a^2)^2} 
\end{align*}

Y por lo tanto

\begin{align*}
    \mathcal{L}^{-1}\left[ \frac{1}{(s^2+a^2)^2}\right](s)  =\frac{1}{2a^3}(\sin(at)-at\cos(t))
\end{align*}

\newpage


\section{Integral de una transformada de Laplace}


Al igual que en la sección anterior, el siguiente resultado puede ser demostrado directamente usando la definición de la transformada de Laplace. \\

\textbf{Proposición: } Sea $f$ una función de orden exponencial, entonces

\begin{align*}
    \int_s^\infty \mathcal{L}[f(t)] (u) du = \mathcal{L}\left[ \frac{f(t)}{t}\right](s)
\end{align*}


Al igual que en la sección anterior, este resultado nos sirve principalmente para calcular transformadas de Laplace de la forma $\displaystyle \mathcal{L}\left[ \frac{f(t)}{t}\right](s)$. \\

\textbf{Ejemplo: } Calculemos la transformada de Laplace de $f(t) =\displaystyle \frac{\sin(at)}{t}$ con $a>0$, sabemos que 


\begin{align*}
    \mathcal{L}\left[ \frac{\sin(at)}{t}\right](s) &= \int_s^\infty \mathcal{L}[\sin(at)] (u) du \\
    &= \int_s^\infty \frac{a}{u^2+a^2}  du \\
    &= \arctan \left( \frac{u}{a}\right) \bigg|_s^\infty \\
    &= \frac{\pi}{2}-\arctan\left( \frac{s}{a}\right)
\end{align*}

\textbf{Ejercicio : } Demuestre que

\begin{align*}
    \mathcal{L}\left[ \frac{\cos(at)}{t}\right](s) 
\end{align*}

No existe, tiende a infinito


\newpage



\section{Resolución de sistemas de ecuaciones diferenciales}

Vamos a ver ahora como aplicar las transformadas de Laplace para resolver sistemas de ecuaciones diferenciales, consideremos

\begin{align*}
   x' &= -6x+2y \\
    y' &= -3x+y \\
\end{align*}

Con las condiciones $y(0)=0$ y $x(0)=1$. Esto es un sistema de ecuaciones diferenciales de $2x2$ en donde $x$ e $y$ son variables de un parámetro $t$. aplicamos transformada de Laplace en ambas igualdades, llamamos $\mathcal{L}[x]=L$ y $\mathcal{L}[y]=M$. luego


\begin{align*}
   sL-1 &= -6L+2M \\
    sM &= -3L+M \\
\end{align*}

Ahora vamos a despejar $L$ y $M$. Despejamos $M$ en la segunda ecuación y obtenemos

\begin{align*}
M &= -\frac{3L}{s-1} \\
\end{align*}

Reemplazando

\begin{align*}
   sL-1 &= -6L-\frac{6L}{s-1} \\
\end{align*}



\begin{align*}
   L &= \frac{s-1}{s^2+5s} \\
   &= \frac{s-1}{s(s+5)}
\end{align*}

y entonces

\begin{align*}
M &= -\frac{3}{s(s+5)} \\
\end{align*}


Usamos fracciones parciales y obtenemos

\begin{align*}
L=\frac{s-1}{s(s+5)} = -\frac{1}{5s}+\frac{6}{5(s+5)}
\end{align*}

\begin{align*}
M= \frac{3}{s(s+5)} = -\frac{3}{5s}+\frac{3}{5(s+5)}
\end{align*}


Luego


\begin{align*}
x(t) &= \mathcal{L}^{-1} \left[ -\frac{1}{5s}+\frac{6}{5(s+5)}\right] \\
&= -\frac{1}{5}+ \frac{6}{5}e^{-5t}
\end{align*}


\begin{align*}
y(t) &= \mathcal{L}^{-1} \left[ -\frac{3}{5s}+\frac{3}{5(s+5)}\right] \\
&= -\frac{3}{5}+ \frac{3}{5}e^{-5t}
\end{align*}

Que es la solución del sistema de ecuaciones diferenciales.









\end{document} 
